% Sección para incluir desde Thesis.tex mediante:
% \input{./chapters/descripcion_problema.tex}

\section{Contexto y motivación del problema}

La Red Metropolitana de Movilidad de Santiago opera diariamente una red de gran
escala compuesta por numerosos recorridos, paradas, terminales y
electroterminales. La incorporación de buses eléctricos permite reducir las
emisiones locales y avanzar hacia una operación más sostenible, pero también
introduce restricciones que no están presentes con la misma intensidad en una
flota convencional. Cada vehículo dispone de una batería de capacidad limitada,
consume energía durante sus desplazamientos y debe acceder oportunamente a
infraestructura de carga para mantener la continuidad del servicio.

La planificación no puede separar las decisiones de transporte de las decisiones
energéticas. Cada corrida programada debe ser atendida por un bus disponible; el
vehículo debe finalizar su actividad anterior, encontrarse o desplazarse hasta
el punto de inicio de la siguiente y disponer de energía suficiente para
completarla. Si necesita recargar, la visita debe insertarse en un intervalo
temporal compatible con la programación y realizarse en un electroterminal
habilitado. A su vez, los desplazamientos sin pasajeros generan tiempo, consumo
de energía y costos que pueden cambiar la conveniencia de una asignación.

La dificultad central consiste en construir cadenas operacionales factibles para
los buses a partir de miles de corridas distribuidas durante el día. Una
solución conveniente al considerar solamente la simultaneidad puede dejar de
ser factible cuando se incorporan la ubicación de los extremos de cada
recorrido, los desplazamientos sin pasajeros, el estado de carga y la
disponibilidad de infraestructura. Del mismo modo, agregar buses reduce la carga
de trabajo promedio por vehículo y puede facilitar la operación, pero aumenta el
costo fijo de flota. El problema requiere evaluar conjuntamente estas
relaciones.

\section{Enunciado, alcance y objetivo}

El problema consiste en determinar un plan diario para la operación de buses
eléctricos que cubra todas las corridas requeridas por la red. El plan debe
decidir cuántos buses utilizar, qué corridas asignar a cada uno, en qué orden
ejecutarlas y cuándo y dónde realizar recargas. Cada corrida debe ser atendida
exactamente una vez y ningún bus puede ejecutar actividades incompatibles en el
tiempo. Además, la secuencia de actividades de cada vehículo debe respetar la
continuidad espacial y mantener su estado de carga dentro de los límites
permitidos.

La cantidad de buses no se fija de manera exógena. Aunque el archivo de
parámetros contiene una referencia de 1.200 buses, este valor no se utiliza como
límite para el dimensionamiento, de acuerdo con la definición posterior de
trabajar con flota irrestricta. En consecuencia, el tamaño de flota pasa a ser
una decisión del problema.

El objetivo general del proyecto es encontrar un plan operacional factible al
menor costo diario posible. Sin embargo, el grupo todavía no ha acordado la
forma matemática definitiva de representar este objetivo ni la manera de
ponderar sus componentes. Entre los costos potencialmente relevantes se
encuentran el uso de buses, los kilómetros recorridos, la energía, las
actividades de carga, los desplazamientos sin pasajeros y los tiempos de espera.
Por esta razón, en esta etapa se declara la minimización de costos como propósito
del proyecto, pero no se propone aún una función objetivo específica.

Las frecuencias entregadas representan el comportamiento ideal esperado de la
red. A partir de ellas se construye una demanda operacional de corridas que el
modelo intentará reproducir al menor costo posible. En consecuencia, el modelo
no podrá reducir costos eliminando salidas o disminuyendo unilateralmente la
frecuencia: el cumplimiento de ese comportamiento ideal será una exigencia de
factibilidad. Se utilizará una misma estructura de modelo, parametrizada para
resolver por separado los patrones laborales, de sábado y de domingo o festivo.
De esta forma, cada tipo de día conservará sus propias corridas, frecuencias y
condiciones de operación, pero los resultados podrán compararse bajo criterios
comunes.

Para esta etapa se supone que todos los buses comienzan con un estado de carga
inicial de 100\%, pueden efectuar recargas parciales y pueden iniciar, recargar
y finalizar su operación en cualquiera de los cinco electroterminales. Estos
supuestos dan flexibilidad al modelo y deberán evaluarse mediante análisis de
sensibilidad, ya que omiten acuerdos operacionales que en la realidad pueden
asociar determinados buses o servicios con terminales específicos.

El entregable esperado es una herramienta computacional reproducible que genere
una asignación factible de corridas y recargas por bus, determine la flota
utilizada y reporte el costo total y sus componentes. La solución también debe
permitir comparar escenarios de demanda, autonomía, precios eléctricos y
capacidad de infraestructura.

\section{Información disponible y filtros aplicados}

La información principal proviene de un conjunto GTFS y de archivos
complementarios sobre vehículos, electroterminales, precios y parámetros. La
Tabla~\ref{tab:fuentes_datos} resume la función de cada fuente dentro del
problema.

\begin{table}[H]
\centering
\caption{Fuentes de datos utilizadas en la descripción del problema}
\label{tab:fuentes_datos}
\small
\begin{tabular}{p{0.29\textwidth}p{0.63\textwidth}}
\toprule
\textbf{Fuente} & \textbf{Información utilizada} \\
\midrule
\texttt{routes.txt} y \texttt{trips.txt} & Agencia, tipo de transporte,
recorrido, patrón, geometría asociada y tipo de día. \\
\texttt{frequencies.txt} & Inicio y término de las ventanas de frecuencia,
intervalo entre salidas y condición \texttt{exact\_times}. \\
\texttt{stop\_times.txt} & Secuencia de paradas y tiempos nominales acumulados
de cada patrón. \\
\texttt{stops.txt}, \texttt{shapes.txt} y
\texttt{shape\_distances.csv} & Localización de paradas, trazado y distancia de
los recorridos. \\
\texttt{vehicles.csv} y \texttt{parameters.csv} & Batería, autonomía,
potencia de carga, límites de SOC y costos unitarios. \\
\texttt{depots.csv} & Identificación, ubicación y capacidad física de los
electroterminales. \\
\texttt{electricity\_prices.csv} & Periodos horarios y precio de la electricidad. \\
\bottomrule
\end{tabular}
\end{table}

El conjunto contiene 427 registros en \texttt{routes.txt}. Para aislar el
problema de buses RED se seleccionaron los registros con
\texttt{agency\_id=RM} y \texttt{route\_type=3}, obteniéndose 417 recorridos.
Este filtro excluye los servicios que no corresponden a buses de la agencia RM,
entre ellos los patrones ferroviarios o de Metro presentes en el GTFS.

Luego se cruzaron los recorridos seleccionados con \texttt{trips.txt} mediante
\texttt{route\_id}. Se conservaron los servicios \texttt{L}, \texttt{S} y
\texttt{D}, interpretados respectivamente como laboral, sábado y domingo o
festivo. El filtro redujo los 26.137 registros originales de
\texttt{trips.txt} a 14.464 patrones de buses: 6.620 laborales, 4.215 de sábado
y 3.629 de domingo o festivo. Cada uno de estos patrones dispone de una fila
correspondiente en \texttt{frequencies.txt}.

Las 14.464 filas seleccionadas en \texttt{frequencies.txt} presentan
\texttt{exact\_times=0}. En consecuencia, el archivo describe un servicio
basado en frecuencias y no una lista de salidas exactas observadas. Para obtener
una instancia operacional determinística se reconstruyeron salidas en
\texttt{start\_time}+n\,\texttt{headway\_secs}, mientras la hora resultante
fuera estrictamente menor que \texttt{end\_time}. Esta convención permite
cuantificar la escala y analizar simultaneidad, pero los horarios obtenidos deben
interpretarse como representativos de la frecuencia declarada.

Para calcular la duración de cada patrón se identificaron su primera y última
parada en \texttt{stop\_times.txt} y se utilizo

\begin{equation}
 T_i = arrival_{i,\mathrm{última}}-departure_{i,\mathrm{primera}}.
 \label{eq:duracion_patron}
\end{equation}

El subconjunto de buses contiene 849.671 registros de paradas. En todos ellos
\texttt{arrival\_time} coincide con \texttt{departure\_time}; por consiguiente,
la base no entrega un tiempo de detención explícito. Esto no hace nula la
duración del recorrido, pues los tiempos acumulados aumentan entre paradas, pero
impide separar tiempo de movimiento y permanencia. Las duraciones empleadas son
nominales y no observaciones de tráfico real.

La distancia se obtuvo desde \texttt{shape\_distances.csv} a través de
\texttt{shape\_id}. Los 14.464 patrones utilizan 839 geometrías distintas y
todas cuentan con una distancia calculada. Finalmente, se estimó el consumo
energético multiplicando la distancia por 1,4 kWh/km. Este factor se deriva de
una batería nominal de 350 kWh y una autonomía declarada de 250 km. Con un SOC
mínimo de 10\%, la energía utilizable por ciclo es 315 kWh y la autonomía
utilizable estimada es 225 km.

\section{Resultados del análisis preliminar}

La expansión de las frecuencias produjo 64.502 corridas laborales, 43.619 de
sábado y 37.377 de domingo o festivo. La Tabla~\ref{tab:escala_demanda} presenta
los principales resultados agregados.

\begin{table}[H]
\centering
\caption{Escala diaria reconstruida por tipo de día}
\label{tab:escala_demanda}
\small
\resizebox{\textwidth}{!}{
\begin{tabular}{lrrrrrr}
\toprule
\textbf{Día} & \textbf{Patrones} & \textbf{Corridas} &
\textbf{Km de servicio} & \textbf{Energía MWh} &
\textbf{$F^{\mathrm{LB}}$} & \textbf{Hora peak} \\
\midrule
Laboral & 6.620 & 64.502 & 1.345.412,69 & 1.883,58 & 6.601 & 08:00 \\
Sábado & 4.215 & 43.619 & 924.133,56 & 1.293,79 & 3.860 & 15:00 \\
Domingo/festivo & 3.629 & 37.377 & 794.791,91 & 1.112,71 & 2.900 & 18:00 \\
\bottomrule
\end{tabular}
}
\end{table}

\begin{table}[H]
\centering
\caption{Características promedio de las corridas reconstruidas}
\label{tab:promedios_corridas}
\small
\begin{tabular}{lrrrr}
\toprule
\textbf{Día} & \textbf{Km/corrida} & \textbf{Min/corrida} &
\textbf{kWh/corrida} & \textbf{Corridas/$F^{\mathrm{LB}}$} \\
\midrule
Laboral & 20,86 & 78,57 & 29,20 & 9,77 \\
Sábado & 21,19 & 72,79 & 29,66 & 11,30 \\
Domingo/festivo & 21,26 & 66,73 & 29,77 & 12,89 \\
\bottomrule
\end{tabular}
\end{table}

Para obtener una referencia mínima de flota, cada corrida reconstruida se
representó mediante el intervalo comprendido entre su salida y su término
nominal. Se reunieron los eventos de inicio y término de toda la red y se
ordenaron cronológicamente. El contador aumenta en uno al comenzar una corrida y
disminuye en uno al finalizarla. Si dos eventos coinciden, se procesa primero el
término, lo que equivale a representar las corridas mediante intervalos
semiabiertos y permitir una reutilización inmediata del bus.

Los máximos del contador fueron 6.601 corridas simultáneas en un día laboral,
3.860 el sábado y 2.900 el domingo o festivo, alcanzados por primera vez a las
08:00, 15:00 y 18:00, respectivamente. Como un bus no puede ejecutar dos
corridas simultáneamente, cualquier solución de cada instancia reconstruida
debe satisfacer

\begin{equation}
 F_L \geq 6\,601, \qquad
 F_S \geq 3\,860, \qquad
 F_D \geq 2\,900.
 \label{eq:cotas_flota_tipo_dia}
\end{equation}

Estas cifras son cotas inferiores condicionadas a la reconstrucción
determinística de las frecuencias. No constituyen todavía flotas factibles, pues
el conteo no asigna corridas a
buses individuales ni considera desplazamientos sin pasajeros, recargas,
compatibilidad geográfica o tiempos de regulación. La sensibilidad a un margen
uniforme ilustra esta limitación: al extender la ocupación de cada corrida en 15
minutos, las cotas aumentan a 7.470 buses en un día laboral, 4.460 el sábado y
3.431 el domingo o festivo, incrementos de 13,2\%, 15,5\% y 18,3\%,
respectivamente. Este escenario no predice atrasos, pero muestra que el resultado
depende de los tiempos de preparación y recuperación incorporados.

En comparación con el día laboral, el sábado tiene 32,4\% menos corridas, 31,3\%
menos kilómetros de servicio y una cota de flota 41,5\% menor. El domingo o
festivo tiene 42,1\% menos corridas, 40,9\% menos kilómetros y una cota de flota
56,1\% menor. La distancia y energía promedio por corrida varían menos de 2\%
entre tipos de día; por ello, la diferencia de flota se explica principalmente
por la distribución temporal y la menor duración promedio de las corridas, y no
por recorridos sustancialmente más cortos.

Desde la perspectiva energética agregada, la flota de la cota laboral dispone al
inicio de $6\,601\times315=2\,079\,315$ kWh utilizables, cifra superior a los
1.883.577,77 kWh requeridos por los kilómetros de servicio. Esto implica una
cota agregada de cero recargas intradía para el escenario laboral, pero no
demuestra que una asignación sin recargas sea factible: la energía se encuentra
distribuida entre buses y algunas cadenas pueden superar la autonomía individual.
Aplicando el mismo cálculo agregado, el sábado presenta un déficit de 77.886,98
kWh y el domingo o festivo uno de 199.208,67 kWh. Si una actividad pudiera
reponer como máximo los 315 kWh utilizables, estos déficits implicarían al menos
248 y 633 cargas equivalentes, respectivamente. Las recargas parciales permitirán
distribuir la energía requerida, aunque potencialmente aumentarán la cantidad de
actividades de carga.

Los resultados también permiten interpretar el intercambio entre flota y
utilización. Al dividir las corridas por su respectiva cota de flota se obtienen
9,77 corridas por bus en un día laboral, 11,30 el sábado y 12,89 el domingo o
festivo. La menor simultaneidad de los fines de semana permitiría, en principio,
reutilizar con mayor intensidad cada bus. Agregar un vehículo reduce solo
marginalmente estos promedios y añade un costo fijo de 250 USD; sin embargo, su
efecto relevante dependerá de si permite evitar desplazamientos vacíos, recargas
costosas o esperas por un valor superior a dicho costo. Esta comparación solo
podrá resolverse al construir cadenas individuales y completar la información
espacial faltante.

\section{Infraestructura y parámetros económicos}

La base identifica cinco electroterminales: Vespucio Norte, El Conquistador, Los
Espinos, La Reina y Santa Rosa. Las capacidades registradas suman 700 buses. El
campo disponible se denomina \texttt{capacity} y describe capacidad física de
albergue; \texttt{depots.csv} no contiene el número de cargadores ni la potencia
del transformador. En consecuencia, los datos actuales no permiten afirmar que
puedan realizarse 700 cargas simultáneas ni imponer completamente la restricción
eléctrica de cada instalación.

La comparación entre los datos revela además una inconsistencia de escala que
debe hacerse explícita. La referencia de 1.200 buses en
\texttt{parameters.csv} y la capacidad agregada de 700 buses en
\texttt{depots.csv} son inferiores a las tres cotas obtenidas al reconstruir
simultáneamente la operación de la agencia RM: 6.601, 3.860 y 2.900 buses. Esto
indica que estos parámetros no pueden imponerse como capacidades duras sobre la
red completa sin una aclaración adicional sobre su alcance. En esta etapa se
descarta el límite de 1.200 buses conforme al supuesto de flota irrestricta, y
la capacidad de los terminales se conserva como información descriptiva, no
como equivalente a capacidad de carga simultánea.

El vehículo de referencia es un BYD K9 con batería de 350 kWh, autonomía de 250
km, potencia máxima de carga de 180 kW y SOC permitido entre 10\% y 100\%. Los
precios eléctricos varían entre 0,10 y 0,22 USD/kWh en seis períodos del día. Los
demás parámetros económicos corresponden a 0,5 USD/km, 0,03 USD por minuto de
espera, 250 USD por bus utilizado y 5 USD por actividad de carga. Estos valores
permiten estimar y comparar costos mientras se define la formulación definitiva
del criterio de optimización. Deben interpretarse como parámetros del caso de
estudio y no necesariamente como costos contractuales observados.

\section{Calidad, limitaciones y necesidades de información}

El análisis permite cuantificar la escala del problema y construir referencias
iniciales, pero aún existen limitaciones que condicionan el modelo final:

\begin{itemize}
 \item Las frecuencias tienen \texttt{exact\_times=0}; las salidas reconstruidas
 son una convención determinística y no horarios exactos observados.
 \item Las duraciones son nominales. No se dispone de realizaciones históricas
 que permitan estimar variabilidad por tráfico, incidentes o comportamiento
 operacional.
 \item Los tiempos de llegada y salida coinciden en todas las paradas de buses,
 por lo que no se representa explícitamente la detención para subida y bajada de
 pasajeros.
 \item No existe una matriz completa de distancia y tiempo sin pasajeros entre
 el final de una corrida y el comienzo de otra, ni entre corridas y
 electroterminales. Sin ella no es posible verificar completamente la
 continuidad espacial.
 \item Los buses no se encuentran individualizados y no se observan sus SOC
 iniciales. El supuesto de inicio al 100\% reemplaza temporalmente esta
 información.
 \item El consumo de 1,4 kWh/km es constante y no incorpora pendiente, velocidad,
 climatización, carga de pasajeros, congestión ni degradación de batería.
 \item La capacidad de albergue de los electroterminales está disponible, pero
 faltan el número de cargadores y la potencia eléctrica máxima de cada terminal.
 \item Permitir que cualquier bus inicie, recargue y termine en cualquier
 electroterminal omite posibles restricciones de propiedad, contrato o
 asignación operacional.
\end{itemize}

Estas limitaciones no impiden formular una primera versión del problema, pero
obligan a distinguir entre factibilidad dentro del modelo y factibilidad real.
Las variables faltantes deben estimarse, obtenerse de una fuente adicional o
tratarse mediante escenarios y análisis de sensibilidad.

\section{Indicadores de desempeño}

La presentación del proyecto define cuatro indicadores preliminares para medir
una solución. Estos cubren las dimensiones de flota, energía, costo y eficiencia
operacional, como se resume en la Tabla~\ref{tab:kpis_proyecto}.

\begin{table}[H]
\centering
\caption{Indicadores de desempeño definidos para evaluar una solución}
\label{tab:kpis_proyecto}
\small
\begin{tabular}{p{0.16\textwidth}p{0.38\textwidth}p{0.22\textwidth}p{0.12\textwidth}}
\toprule
\textbf{Dimensión} & \textbf{Indicador} & \textbf{Valor de referencia} &
\textbf{Importancia} \\
\midrule
Flota & Buses utilizados & 6.601 (L); 3.860 (S); 2.900 (D) & 3/5 \\
Energía & Porcentaje de ocupación de terminales & 90\% (630 buses) & 2/5 \\
Costo & Costo total diario & 2.545.946 USD & 4/5 \\
Eficiencia & Porcentaje del tiempo del bus fuera de servicio & Menor a 20\% & 2/5 \\
\bottomrule
\end{tabular}
\end{table}

Los valores indicados son referencias preliminares y no resultados óptimos. Las
cifras de flota corresponden a las cotas de simultaneidad de las instancias
laboral, sábado y domingo o festivo reconstruidas. La ocupación de 90\% equivale
a 630 de los 700 espacios
físicos nominales registrados en los electroterminales y no debe interpretarse
como capacidad de carga simultánea. El costo total diario de referencia fue
calculado para el escenario laboral y se utilizará para comparar alternativas
cuando se acuerde su composición definitiva. Finalmente,
el indicador de eficiencia busca que el tiempo asociado a espera, recarga y
desplazamientos sin pasajeros represente menos de 20\% del tiempo total de
disponibilidad del bus. De manera complementaria, toda solución deberá mantener
una cobertura de 100\% de las corridas del comportamiento ideal reconstruido.
